\documentclass[12pt,a4paper]{article}
\usepackage{fontspec}
\setmainfont{Times New Roman}
\usepackage[margin=1in]{geometry}   % 设置边距
\usepackage{graphicx}               % 插入图片
\usepackage{setspace}               % 控制行距
\setlength{\parindent}{0pt}         % 段落不缩进
\usepackage[hyphens]{url}
\usepackage[colorlinks=true, linkcolor=black, urlcolor=blue, citecolor=black]{hyperref}
\hypersetup{breaklinks=true}
\usepackage[strings]{underscore}
\makeatletter
\renewcommand{\section}{\@startsection{section}{1}{\z@}%
  {-3.5ex \@plus -1ex \@minus -.2ex}%
  {2.3ex \@plus.2ex}%
  {\normalfont\normalsize\bfseries}}
\makeatother

\begin{document}


%-------------------------
% Title and Contact
%-------------------------
\begin{center}
    {\textbf{Xufeng Duan}}\\
    \textbf{Email: } xufengduan@cuhk.edu.hk \textbf{Tel: }66731863 \quad  \\
\textbf{Research interests:} Psycholinguistics, Neurolinguistics, Mechanistic Interpretability
\end{center}

%-------------------------
% Education
%-------------------------
\vspace{-2em} 
\section*{EDUCATION}
\vspace{-.5em}
\hrule
\vspace{.2em}
\textbf{The Chinese University of Hong Kong} \hfill 2020--2025\\
Department of Linguistics and Modern Languages \hfill \textit{Ph.D.} in Linguistics

% \textbf{The Chinese University of Hong Kong} \hfill 2019--2020\\
% \href{https://cuhklpl.github.io/}{Department of Linguistics and Modern Languages} \hfill \textit{M.A.} in Linguistics

\textbf{The Chinese University of Hong Kong} \hfill 2019--2020\\
Department of Linguistics and Modern Languages \hfill \textit{M.A.} in Linguistics

\textbf{Sun Yat-sen University} \hfill 2015--2019\\
Neurolinguistics Laboratory \hfill \textit{B.A.} in Chinese
\vspace{\baselineskip}

%-------------------------
% Working Experience
%-------------------------
\vspace{-2em} 
\section*{CURRENT POSITION}
\vspace{-.5em}
\hrule
\vspace{.2em}
\textbf{The Chinese University of Hong Kong} \hfill 2025--Present\\
Brain and Mind Institute \hfill \textit{Research Assistant Professor}
\vspace{\baselineskip}


\vspace{-2em}
\section*{PREVIOUS RELEVANT RESEARCH WORK}
\vspace{-.5em}
\hrule
\vspace{.2em}
My research integrates \textbf{artificial intelligence (AI) and psycholinguistics} to develop an emerging field of \textit{AI for Psychology}. This proposal extends a coherent research trajectory in which I have progressively established the three methods of the current project: multi-agent reasoning, automated stimulus construction, and AI-based behavioural simulation.\\
\hspace*{2em}First, I developed \textbf{multi-agent frameworks} that emulate scientific reasoning. \textit{StimMAS} (Duan, Xiao, Kan, \& Cai, 2025) introduced a multi-agent system that generates, critiques, and validates linguistic stimuli, demonstrating how distributed AI agents can coordinate hypothesis testing under psycholinguistic constraints. Related work on automated literature review and collaborative generation (\textit{Tang, Duan, \& Cai, 2025, EMNLP}) and applied multi-agent tools (\textit{Zhao \& Duan, 2025}) further establish the feasibility of agent-based hypothesis generation.\\
\hspace*{2em}Second, I have shown that AI can \textbf{construct experimental stimuli} comparable to human-designed materials. Using \textit{StimMAS}, we replicated classic psycholinguistic effects with 300 participants across three paradigms, confirming the reliability of AI-generated stimuli. \\
\hspace*{2em}Third, our team pioneered the use of \textbf{LLMs as computational participants}. The \textit{MacBehaviour} framework (Duan, Li, \& Cai, 2025, \textit{Behavior Research Methods}) formalized AI-based behavioural experimentation, revealing human-like sensitivity to mental lexicon, syntax, semantics, and grammaticality (\textit{Cai, Duan, Haslett, Wang, \& Pickering., 2024, cited 142 times as of 17/10/2025; Xiao, Duan, Haslett \& Cai., 2025; Duan, Xiao, Tang \& Cai., 2025;  Qiu, Duan \& Cai., 2025}). Eye-tracking comparisons (\textit{*Wang, *Duan, \& Cai, 2025}) further demonstrated strong alignment between human gaze patterns and model attention, as well as between ERP amplitude and LLM surprisal(\textit{Wu, Duan, \& Cai, 2025}). Complementary work on model interpretability (\textit{Duan et al., 2025; Zhou et al., 2025, COLING}) links these methods to theoretical on language processing. \\
\hspace*{2em}Collectively, this body of research has produced validated tools (\textit{StimMAS}, \textit{MacBehaviour}), datasets, and interdisciplinary methods that render the proposed GRF project both technically feasible and conceptually mature.

\vspace{\baselineskip}
%-------------------------
% Publications
%-------------------------
\vspace{-2em} 
\section*{PUBLICATION}
\vspace{-0.5em}
\hrule
\vspace{.2em}

\underline{Five most representative publications in recent five years (* corresponding author)}
\vspace{1em}

\hangindent=1.5em \hangafter=1 \textbf{Duan, X.}, Li, S., \& Cai, Z. G. (2025). MacBehaviour: An R package for behavioural experimentation on large language models. \textit{Behavior Research Methods}, 57(1), 19. 

\hangindent=1.5em \hangafter=1 \textbf{Duan, X.}, Zhou, X., Xiao, B., \& Cai, Z. (2025). Unveiling language competence neurons: A psycholinguistic approach to model interpretability. In O. Rambow, L. Wanner, M. Apidianaki, H. Al-Khalifa, B. D. Eugenio, \& S. Schockaert (Eds.), \textit{Proceedings of the 31st International Conference on Computational Linguistics (COLING)} (pp. 10148–10157). Association for Computational Linguistics.

\hangindent=1.5em \hangafter=1 \textbf{Duan, X.}, Xiao, B., Kan, B., \& Cai, Z. G. (2025). StimMAS: A Multi-Agent Framework for Automated Linguistic Stimulus Construction in Psychological Research. \textit{PsyArXiv preprint, under review at AMPPS}. 

\hangindent=1.5em \hangafter=1 Tang, X., \textbf{Duan, X.}$^{*}$ , \& Cai, Z. G. (2025). Large Language Models for Automated Literature Review: An Evaluation of Reference Generation, Abstract Writing, and Review Composition. \textit{In Proceedings of the 2025 Conference on Empirical Methods in Natural Language Processing (EMNLP), Main Conference}. 


\hangindent=1.5em \hangafter=1 Wu, H., \textbf{Duan, X.}, \& Cai, Z. (2025). Distinct social-linguistic processing between humans and large audio-language models: Evidence from model-brain alignment. In T. Kuribayashi, G. Rambelli, E. Takmaz, P. Wicke, J. Li, \& B.-D. Oh (Eds.), \textit{Proceedings of the Workshop on Cognitive Modeling and Computational Linguistics (CMCL)} (pp. 135–143). NAACL 2025. 

\vspace{\baselineskip}
\underline{Five representative publications beyond the recent five-year period }
\vspace{.5em}\\
N.A.\\

\underline{Related publications (${\dagger}$ co-first author)}
\vspace{.5em}

\hangindent=1.5em \hangafter=1 Cai, Z., \textbf{Duan, X.}, Haslett, D., Wang, S., \& Pickering, M. (2024). Do large language models resemble humans in language use? \textit{Proceedings of the Workshop on Cognitive Modeling and Computational Linguistics (CMCL)} (pp. 37–56).

\hangindent=1.5em \hangafter=1 Wang, S.$^{\dagger}$, \textbf{Duan, X.$^{\dagger}$}, \& Cai, Z. (2024). A Multimodal Large Language Model “Foresees” Objects Based on Verb Information but Not Gender. \textit{Proceedings of the 28th Conference on Computational Natural Language Learning (CONLL)} (pp. 435–441).

\hangindent=1.5em \hangafter=1 Wu, H., \textbf{Duan, X.}, \& Cai, Z. G. (2024). Speaker Demographics Modulate Listeners' Neural Correlates of Spoken Word Processing. \textit{Journal of Cognitive Neuroscience}, 36(10), 2208–2226. 

\hangindent=1.5em \hangafter=1 \textbf{Duan, X.}, \& Cai, Z. G. (2024). Chinese Character Processing. \textit{Oxford Research Encyclopedia of Linguistics}. Oxford University Press. 

\hangindent=1.5em \hangafter=1 Zhou, X., Chen, D., Cahyawijaya, S., \textbf{Duan, X.}, \& Cai, Z. G. (2025). Linguistic Minimal Pairs Elicit Linguistic Similarity in Large Language Models. \textit{Proceedings of the 31st International Conference on Computational Linguistics (COLING)} (pp. 6866–6888). Association for Computational Linguistics.

\hangindent=1.5em \hangafter=1 \textbf{Duan, X.}, Xiao, B., Tang, X., \& Cai, Z. G. (2024). HLB: Benchmarking LLMs' Humanlikeness in Language Use. \textit{arXiv:2409.15890 [cs.CL]}.

\hangindent=1.5em \hangafter=1 Zhao, N., \& \textbf{Duan, X.} (2025). Enhancing Conference Interpreting with Computer-assisted Interpreting Tools: A Multi-agent System. \textit{PsyArXiv preprint}.





\vspace{\baselineskip}

%-------------------------
% Research Projects
%-------------------------
\vspace{-2em} 
\section*{RESEARCH PROJECTS}
\vspace{-.5em}
\hrule
\vspace{.2em}

\begin{itemize}
\item Participate, General Research Fund 2021/22 (14613722), \textit{Neurocognitive mechanisms of character amnesia in Chinese handwriting}, RGC, HK\$803,050, 01/01/2023 – 31/12/2025. PI: Zhenguang Cai (CUHK).

    
\end{itemize}
\end{document}


